% ---------------------------------------------------------------------
\section{Evidence Log: Kill-Test on Real Hardware}
\label{sec:evidence}
% ---------------------------------------------------------------------

This chapter records the de-risking experiment run \emph{before} any HyperPot
code was written, to verify the core premise on the actual demo laptops rather
than relying only on documentation.

\subsection{Setup}
\begin{table}[h]
\centering\footnotesize
\caption{Kill-test configuration (26 Oct 2026)}
\label{tab:killtest}
\begin{tabularx}{\linewidth}{L{0.26\linewidth}Y}
\toprule
Item & Value \\
\midrule
Laptop A (main process) & ROG, RTX 3050 Laptop GPU (3767 MiB), Manjaro Linux, CUDA 13.3 \\
Laptop B (remote node) & TUF, RTX 4050, Windows, llama.cpp b11900 CPU build \\
Engine & llama.cpp v0.7.0-dev (commit caa8191), built with \code{-DGGML\_CUDA=ON -DGGML\_RPC=ON} on the ROG \\
Model & Llama-3.2-1B-Instruct, Q4\_K\_M GGUF (771 MB) \\
Link & Wi-Fi; B exposed \code{ggml-rpc-server -H 0.0.0.0 -p 50052} \\
Connectivity & TCP port 50052 reachable from A (verified) \\
\bottomrule
\end{tabularx}
\end{table}

\subsection{Run 1 --- normal distributed inference (baseline)}
A full generation ran across A~+~B over RPC and \textbf{completed successfully},
producing a coherent multi-paragraph story. Final stats:
\code{[ Prompt: 52.7 t/s | Generation: 7.7 t/s ]}, then \code{Exiting...}.
\textbf{Conclusion:} distributed inference across two heterogeneous laptops over
Wi-Fi \emph{works}. \VER

\subsection{Run 2 --- remote node lost (the premise)}
A longer run was started. During model data transfer to the remote node, the
Wi-Fi link to B dropped packets and the remote node became unreachable. The main
process on A did \textbf{not} recover --- it aborted the entire process:
\begin{lstlisting}[style=plain]
ggml/src/ggml-rpc/ggml-rpc.cpp:685: Remote RPC server crashed
    or returned malformed response
E send failed (bytes_sent=6877944, size_to_send=2559545)
Aborted (core dumped)
\end{lstlisting}

\subsection{What this proves}
\begin{itemize}[leftmargin=1.4em,nosep]
  \item \VER\ \textbf{llama.cpp RPC has no fault tolerance on this exact
  hardware.} Loss of the remote node terminates the whole process
  (\code{ggml-rpc.cpp:685} $\rightarrow$ \code{Aborted}), matching the
  documentation (\S\ref{sec:landscape}).
  \item \VER\ \textbf{The failure is realistic, not contrived.} It was triggered
  by ordinary Wi-Fi instability --- no deliberate kill was even needed. This is
  precisely the scenario HyperPot targets.
  \item \textbf{Answer to ``does it restart from the beginning?''} No --- plain
  llama.cpp does not restart at all; it crashes and must be re-launched by hand.
  HyperPot's contribution is to make the controller detect the loss, re-route the
  work, and \textbf{automatically re-run the request} so generation continues.
\end{itemize}

\subsection{Note for the live demo}
The crash here came partly from random Wi-Fi drops during a large transfer to a
CPU-only node. For a reproducible stage demo, the ``kill'' should be a
\textbf{deliberate, clean action} (Ctrl+C on the remote server, or lid-close)
rather than relying on chance packet loss. The before/after contrast ---
llama.cpp aborts vs HyperPot recovers --- is the core of the demonstration. Both
log files (\code{killtest.log}, \code{killtest2.log}) are retained as evidence.
